
\documentclass{article}
\usepackage{colortbl}
\usepackage{makecell}
\usepackage{multirow}
\usepackage{supertabular}

\begin{document}

\newcounter{utterance}

\centering \large Interaction Transcript for game `wordle\_withclue', experiment `medium\_frequency\_words\_clue\_no\_critic', episode 16 with selfrefine{-}Qwen3.5{-}9B.
\vspace{24pt}

{ \footnotesize  \setcounter{utterance}{1}
\setlength{\tabcolsep}{0pt}
\begin{supertabular}{c@{$\;$}|p{.15\linewidth}@{}p{.15\linewidth}p{.15\linewidth}p{.15\linewidth}p{.15\linewidth}p{.15\linewidth}}
    \# & \multicolumn{2}{c}{Player} && \multicolumn{2}{c}{Game Master} \\
    \hline

    \theutterance \stepcounter{utterance}  
    & & & \multicolumn{4}{p{0.6\linewidth}}{
        \cellcolor[rgb]{0.9,0.9,0.9}{
            \makecell[{{p{\linewidth}}}]{
                \texttt{\tiny{[P1$\langle$GM]}}
                \texttt{You are a language wizard who likes to guess words by using the given rules.} \\
\\ 
\texttt{Welcome to Wordle! You have six attempts to guess the target word, a valid English word of five lowercase letters (a{-}z). Please use the tags "explanation:" and "guess:" to provide a concise explanation for each guess.} \\
\\ 
\texttt{To help you make an informed guess, you will receive a clue for the word, such as} \\
\texttt{clue: snowy white.} \\
\\ 
\texttt{Here is an example guess based on the clue:} \\
\texttt{explanation: In the fairy tail Snow White, the girl is killed because she eats a poisoned apple. And the word apple has 5 letters.} \\
\texttt{guess: apple} \\
\\ 
\texttt{After each guess, your answer will be validated, and you will receive feedback indicating which letters are correct (green), which letters are correct but in the wrong position (yellow), and which letters are incorrect (red). This feedback can be useful in determining which letters to include or exclude in your next guess.} \\
\\ 
\texttt{For example, the feedback for "apple" might be:} \\
\texttt{guess\_feedback: a<yellow> p<yellow> p<green> l<yellow> e<red>} \\
\\ 
\texttt{The explanation should contain details about how both the guess\_feedback and clue work together in tandem to formulate a new guess.} \\
\\ 
\texttt{Do not generate any other text.} \\
\\ 
\texttt{Let's begin with your first guess.} \\
\\ 
\texttt{clue: shrewd} \\
            }
        }
    }
     \\ \\

    \theutterance \stepcounter{utterance}  
    & & \multicolumn{4}{p{0.6\linewidth}}{
        \cellcolor[rgb]{0.95,0.95,0.95}{
            \makecell[{{p{\linewidth}}}]{
                \texttt{\tiny{[P1$\leftrightarrow$P]}}
                \texttt{{[}DRAFT{]}} \\
\texttt{explanation: The clue "shrewd" suggests the target word is related to being clever or cunning. A common synonym for shrewd is "sharp," which fits the five{-}letter requirement. Additionally, "sharp" contains common letters like 's', 'h', 'a', 'r', and 'p' that might appear in the target word, making it a strategic first guess to gather information.} \\
\texttt{guess: sharp} \\
            }
        }
    }
    & \\ \\

    \theutterance \stepcounter{utterance}  
    & & \multicolumn{4}{p{0.6\linewidth}}{
        \cellcolor[rgb]{0.95,0.95,0.95}{
            \makecell[{{p{\linewidth}}}]{
                \texttt{\tiny{[P1$\leftrightarrow$P]}}
                \texttt{{[}SELF{-}FEEDBACK{]}} \\
\texttt{FORMAT: ok} \\
\texttt{RULES: ok} \\
\texttt{STRATEGY: problem} \\
\texttt{REASON: The guess "sharp" is a valid 5{-}letter word, but it is a poor strategic choice for a first guess because it contains many low{-}frequency letters (like 'h' and 'p') and lacks the high{-}frequency vowels and consonants needed to maximize information gain.} \\
            }
        }
    }
    & \\ \\

    \theutterance \stepcounter{utterance}  
    & \multicolumn{4}{p{0.6\linewidth}}{
        \cellcolor[rgb]{0.561,0.737,0.561}{
            \makecell[{{p{\linewidth}}}]{
                \texttt{\tiny{[P1$\rangle$GM]}}
                \texttt{explanation: The clue "shrewd" suggests a word meaning clever or cunning. While "sharp" is a valid synonym, it is a poor first guess because it lacks high{-}frequency vowels (like 'e', 'a', 'o') and contains low{-}frequency letters ('h', 'p') that yield less information. A better strategy for the first guess is to use a word with many common letters to maximize feedback. "crane" is a strong alternative as it contains 'r', 'n', 'e', 'a' which are very common, helping to quickly identify or eliminate key letters.} \\
\texttt{guess: crane} \\
            }
        }
    }
    & & \\ \\

    \theutterance \stepcounter{utterance}  
    & & & \multicolumn{2}{p{0.3\linewidth}}{
        \cellcolor[rgb]{0.9,0.9,0.9}{
            \makecell[{{p{\linewidth}}}]{
                \texttt{\tiny{[GM$|$GM]}}
                \texttt{attempts = 1} \\
\texttt{target\_word = canny} \\
\texttt{target\_word\_clue = shrewd} \\
\texttt{guess = crane} \\
\texttt{guess\_feedback = c<green> r<red> a<yellow> n<green> e<red>} \\
            }
        }
    }
    & & \\ \\

    \theutterance \stepcounter{utterance}  
    & & & \multicolumn{4}{p{0.6\linewidth}}{
        \cellcolor[rgb]{0.9,0.9,0.9}{
            \makecell[{{p{\linewidth}}}]{
                \texttt{\tiny{[P1$\langle$GM]}}
                \texttt{guess\_feedback: c<green> r<red> a<yellow> n<green> e<red>} \\
\\ 
\texttt{Provide your response only in this format.} \\
\texttt{explanation: details} \\
\texttt{guess: word} \\
            }
        }
    }
     \\ \\

    \theutterance \stepcounter{utterance}  
    & & \multicolumn{4}{p{0.6\linewidth}}{
        \cellcolor[rgb]{0.95,0.95,0.95}{
            \makecell[{{p{\linewidth}}}]{
                \texttt{\tiny{[P1$\leftrightarrow$P]}}
                \texttt{{[}DRAFT{]}} \\
\texttt{explanation: The feedback for "crane" reveals that 'c' and 'n' are correct and in the right positions (green), while 'a' is correct but in the wrong position (yellow), and 'r' and 'e' are incorrect (red). Since the clue is "shrewd," which implies cleverness or sharpness, and we know the word starts with 'c' and has 'n' as the fourth letter, we need a word fitting the pattern c\_n\_ with 'a' not in the third position. The word "canon" fits this pattern (c{-}a{-}n{-}o{-}n) but has 'a' in the second spot, which contradicts the yellow 'a' from "crane" (where 'a' was third). Wait, in "crane", 'a' was the third letter and got yellow, meaning 'a' is in the word but not third. In "canon", 'a' is second. Let's re{-}evaluate. The pattern is c \_ \_ n \_. 'a' is present but not at index 2. 'r' and 'e' are out. The clue "shrewd" might hint at "cunning" but that's 7 letters. How about "canyon"? No, 6 letters. Let's try "cabin"? No 'a' yellow issue. What about "canal"? c{-}a{-}n{-}a{-}l. 'a' is at 2 and 4. In "crane", 'a' was at 3. So 'a' is not at 3. In "canal", 'a' is at 2 and 4. This is possible. But "shrewd" doesn't strongly suggest "canal". Maybe "crown"? No 'a'. "coven"? No 'a'. "coven" has 'o', 'v', 'e' (e is red). "coven" doesn't have 'a'. The word must have 'a'. "cabin" has 'a' at 2. "cabin" {-}> c{-}a{-}b{-}i{-}n. 'a' is at 2. In "crane", 'a' was at 3. So 'a' is not at 3. This is consistent. But "cabin" doesn't mean shrewd. "canny" means shrewd or astute. c{-}a{-}n{-}n{-}y. 'a' is at 2. 'n' is at 3 and 4. In "crane", 'n' was at 4 and got green. So 'n' is at 4. In "canny", 'n' is at 3 and 4. This fits. 'a' is at 2, which is not 3, so yellow in "crane" is consistent. 'r' and 'e' are} \\
            }
        }
    }
    & \\ \\

    \theutterance \stepcounter{utterance}  
    & & \multicolumn{4}{p{0.6\linewidth}}{
        \cellcolor[rgb]{0.95,0.95,0.95}{
            \makecell[{{p{\linewidth}}}]{
                \texttt{\tiny{[P1$\leftrightarrow$P]}}
                \texttt{{[}SELF{-}FEEDBACK{]}} \\
\texttt{FORMAT: problem} \\
\texttt{RULES: problem} \\
\texttt{STRATEGY: problem} \\
\texttt{REASON: The response fails the format requirement by not using the specified tags, violates game rules by proposing a 6{-}letter word ("canny" is 5 letters but the draft discusses "canyon" and "canon" which are invalid lengths or patterns, and the draft itself is cut off), and lacks a clear, valid strategy due to the incomplete and confused reasoning.} \\
            }
        }
    }
    & \\ \\

    \theutterance \stepcounter{utterance}  
    & \multicolumn{4}{p{0.6\linewidth}}{
        \cellcolor[rgb]{0.561,0.737,0.561}{
            \makecell[{{p{\linewidth}}}]{
                \texttt{\tiny{[P1$\rangle$GM]}}
                \texttt{explanation: The feedback from "crane" confirms 'c' (1st) and 'n' (4th) are correct, while 'a' is present but not in the 3rd position. Letters 'r' and 'e' are incorrect. The clue "shrewd" suggests a synonym for clever or cunning. The word "canny" fits the pattern c\_a\_n\_n\_y, placing 'a' in the 2nd position (consistent with the yellow feedback) and 'n' in the 4th position (consistent with the green feedback). It also aligns perfectly with the meaning of the clue.} \\
\texttt{guess: canny} \\
            }
        }
    }
    & & \\ \\

    \theutterance \stepcounter{utterance}  
    & & & \multicolumn{2}{p{0.3\linewidth}}{
        \cellcolor[rgb]{0.9,0.9,0.9}{
            \makecell[{{p{\linewidth}}}]{
                \texttt{\tiny{[GM$|$GM]}}
                \texttt{attempts = 2} \\
\texttt{target\_word = canny} \\
\texttt{target\_word\_clue = shrewd} \\
\texttt{guess = canny} \\
\texttt{guess\_feedback = c<green> a<green> n<green> n<green> y<green>} \\
            }
        }
    }
    & & \\ \\

    \theutterance \stepcounter{utterance}  
    & & & \multicolumn{2}{p{0.3\linewidth}}{
        \cellcolor[rgb]{0.9,0.9,0.9}{
            \makecell[{{p{\linewidth}}}]{
                \texttt{\tiny{[GM$|$GM]}}
                \texttt{game\_result = WIN} \\
            }
        }
    }
    & & \\ \\

\end{supertabular}
}

\end{document}
